\section{Related literature}
\label{sec:literature}

This paper sits at the intersection of four literatures: the empirical incidence of energy price shocks on households; the distinction between vertical and horizontal redistribution in energy policy design; the methodological literature on what a static consumption-side microsimulation can and cannot recover; and the macroeconomics of geopolitical risk and oil supply. I take them in that order, and close by naming the gap the paper occupies.

\subsection{The unresolved question: compensate incomes, or subsidise prices?}

The most useful way to enter this literature is through a disagreement that is rarely stated as one, because the two sides publish in different venues. On one side stands what has become the grey-literature consensus of the 2022--23 European energy crisis. \citet{ari2022}, in the IMF working paper that set the terms of the debate, argue that governments should let retail prices rise to their market level and compensate households through income transfers, on the standard ground that a price subsidy destroys the conservation incentive at precisely the moment when supply is scarce and thereby buys distributional relief at an unnecessary efficiency cost. The Bruegel national fiscal policy tracker supplied the empirical indictment: only around a third of European support during the crisis was targeted, the remainder being untargeted price measures \citep{sgaravatti2023}. \citet{amores2023} reach compatible conclusions in EUROMOD for the EU as a whole, and the European Central Bank's own review took the same line \citep{ecb2023}. The slogan---compensate incomes, never prices---became the default recommendation of virtually every international institution, and it is the intellectual backdrop against which the Chancellor's rejection of a repeated Energy Price Guarantee should be read.

On the other side stands the peer-reviewed frontier. \citet{levell2025}, in work conditionally accepted at the \emph{Journal of Political Economy}, solve for optimal energy-crisis policy in a framework estimated on UK data and reach a conclusion the consensus does not admit. Absent any policy, the mean welfare loss from the 2022--23 shock was around six per cent of income; the package the UK government actually implemented reduced both the mean and the dispersion of those losses at an efficiency cost of twelve per cent of its revenue cost---which is a real cost, but a modest one relative to the redistribution achieved. Crucially, their optimal policy includes a \emph{strictly positive price subsidy}. The subsidy shrinks as the transfer instrument becomes better able to condition on both income and past energy usage, but it does not vanish, because energy usage carries information about household need that income alone does not. This is not a rejection of the consensus so much as a demonstration of what the consensus assumes: that a sufficiently well-targeted income transfer is available. Where the transfer instrument is coarse---and the British instrument is coarse, since roughly 40 per cent of households struggling to heat their homes are not on means-tested benefits---the case for a price instrument survives the theory. The UK's live choice between a means-tested social tariff and a universal discounted block is, in effect, an empirical instance of exactly this question, and it is being decided without the corresponding UK evidence.

\subsection{Incidence: pounds versus per cent}

The empirical incidence literature contains two facts that are usually reported separately and that point in opposite directions. \citet{fetzer2024}, in the paper that is central to this paper's headline figure, match 22.2 million Energy Performance Certificates to National Energy Efficiency Data-Framework meter records and establish that exposure in \emph{cash} terms rises with affluence: more affluent areas occupy larger, older and less thermally efficient dwellings, consume more energy, and therefore lose more pounds when the unit price rises. The direct policy implication, which they draw, is that the Energy Price Guarantee---a per-unit subsidy---disproportionately benefited the households least in need of it, and they propose a two-tier tariff in which an initial block is priced below the marginal block. They are explicit that this is the only distributional statement their data support: lacking income measures, they cannot express energy expenditure as a share of income, so their result speaks to the cash arm alone and neither confirms nor contradicts the percentage arm. The countervailing fact is the older and less contested one: because energy is a necessity whose Engel curve falls with income, the same shock costs poorer households far more as a \emph{share} of resources. \citet{advani2013} document this for the UK in the pre-crisis period, reporting an energy budget share of 15.8 per cent of expenditure in the bottom decile against 3.3 per cent in the top. Inflation itself therefore differs systematically across the distribution, although the sharpest available UK measurement of that divergence is not an energy result: using household scanner data on fast-moving consumer goods, \citet{chen2024} find that grocery-price inflation cumulated over 2021Q3--2023Q3 was 7.7 percentage points lower for the top expenditure decile than for the bottom, driven by ``cheapflation''---a systematic rise in the relative price of cheaper product varieties. That figure is cumulative, and it covers groceries rather than energy; it establishes that the mechanism operates, not its magnitude in the energy basket. The ONS Household Costs Indices now publish the household-group price series that make the divergence visible in official statistics.

Both facts are correct, and a distributional analysis that reports only one of them will support the wrong policy. A per-unit subsidy defended on the grounds that poor households lose most in percentage terms will in fact spend most of its money at the top; a strictly means-tested transfer defended on the grounds that affluent households lose most in cash will fail to reach the fuel-poor households outside the benefit system. What has been missing is a presentation that shows both at once, on the same households. That is what the paired pounds-versus-per-cent profile in Section~\ref{sec:results} supplies.

\subsection{Vertical and horizontal redistribution}

A closely related literature argues that the decile table---the standard output of distributional analysis, including PolicyEngine's---systematically understates the redistribution that energy policy actually performs. \citet{cronin2019} show for US carbon pricing that horizontal redistribution \emph{within} income deciles exceeds vertical redistribution \emph{between} them: two households at the same income can face wholly different exposure depending on dwelling type, heating fuel, location and household composition, and the between-decile average conceals this entirely. \citet{douenne2020} makes the same point for France with a sharper policy edge: a flat recycling of carbon tax revenue is progressive on average, yet leaves a large mass of horizontal losers concentrated among rural and off-gas-grid poor households, precisely the group least visible in the aggregate. \emph{No UK counterpart to Douenne exists.} This is a genuine gap, and I state it as one: Britain has a large off-gas-grid population, an unusually inefficient housing stock, and a heating-fuel geography that maps imperfectly onto income, so the French result should be expected to travel, but nobody has established it.

\citet{sallee2019} supplies the theoretical bound on how far compensation can go. When transfers must be conditioned on observable characteristics---income, benefit receipt, household size---rather than on the unobserved heterogeneity that actually determines exposure, compensating \emph{all} losers is infeasible: any feasible scheme leaves some genuine losers uncompensated, and pushing the uncompensated share towards zero requires overcompensating a large number of non-losers. This converts an abstract worry into a reportable statistic. The right question to ask of a social tariff is not what it delivers to the bottom decile on average, but what share of losing households \emph{within} each decile it leaves whole---and a household-level decomposition of winners and losers within each decile produces exactly that object. It is the number the UK debate has not been given, and Section~\ref{sec:policy} reports it for each of the five live proposals.

\subsection{What a static microsimulation cannot see}

The methodological literature imposes discipline on the claims a paper like this may make. \citet{kanzig2023} is the principal threat. Using identified energy-price shocks and US household data, he finds that poorer households suffer disproportionately not mainly through their energy budget shares but through income and employment, and that general-equilibrium and indirect effects account for roughly two-thirds of the total household response. A static consumption-side microsimulation, which holds employment, wages and prices of non-energy goods fixed, therefore recovers at most a third of true incidence. This is not a reason to abandon the exercise---the use-side loss is the object that compensation policy is designed to offset, and it is the object the Treasury will be asked to fund---but it does bound the interpretation, and I treat every estimate here as a lower bound on total welfare incidence.

\citet{goulder2019} decompose the mechanism. Use-side effects of energy pricing are regressive, because energy budget shares fall with income; source-side effects are progressive, because the incidence falls partly on capital and on the returns to energy-intensive production, which are concentrated at the top. Microsimulation of household consumption sees only the use side. \citet{ohlendorf2021}, in a meta-analysis of the distributional literature, show that whether an energy price change is judged progressive or regressive depends systematically on methodological choices---income versus expenditure as the resource measure, the treatment of the source side, the country, and the time horizon---which is a warning against reporting a single headline sign. \citet{pallotti2024} deliver the most disruptive finding of the four: in euro-area data covering the 2021--23 inflation surge, most of the cross-household heterogeneity in welfare losses was driven by nominal net asset positions---who held nominal debt and who held nominal savings---rather than by differences in the energy consumption basket at all. PolicyEngine UK carries wealth imputed from the Wealth and Assets Survey, so the ingredients for a British version of that decomposition exist, but I do not attempt it here.

Two further methodological anchors are conventional. The first-order approximation used in Section~\ref{sec:methodology} is Deaton's: the welfare cost of a small price change is the pre-shock quantity times the price change, which is an upper bound on the compensating variation once substitution is allowed. Where demand systems are required, the standard references are \citet{deaton1980} for the Almost Ideal Demand System and \citet{banks1997} for its quadratic extension; \citet{odonoghue2023} set out the PRICES framework for combining price shocks with tax-benefit microsimulation---in a Pakistani application, but the procedure is general---and it is the closest procedural precedent for what I do. The relevant British limitation is that PolicyEngine UK has no elasticity or pass-through block, so the first-order approximation is not a choice among alternatives but the only estimator the model supports.

\subsection{Geopolitical risk, oil supply and the macro side}

The macro strand supplies the shock. \citet{caldara2022} construct the Geopolitical Risk index from newspaper text and establish it as the standard measure, with threat and act subindices published daily. \citet{caldara2025} extend the analysis to the transmission mechanism and find that geopolitical risk \emph{raises} inflation while lowering activity---a genuinely stagflationary configuration---with commodity-price and currency channels dominating, which is the pattern the UK experienced after February 2026. \citet{verduzco2026} sharpen the identification by isolating geopolitical events that specifically reduce oil supply, and show that oil \emph{inventories} are the key margin along which such shocks propagate, a result that matters here because inventory buffers determine how much of a Hormuz disruption reaches the retail price at all.

The most useful complication comes from \citet{boe2025}, who separate geopolitical ``Acts'' from ``Threats'' and find that Acts---realised military events---produce \emph{negative} oil price and inflation responses on average, since markets have typically priced the threat in advance and the act resolves uncertainty. This directly contradicts the naive stagflation prior and serves as a foil: any claim that the 2026 episode was inflationary must be defended against the possibility that it was an Act rather than a Threat, and I address the point when calibrating the scenario. The structural foundations of the oil-market literature are \citet{hamilton1983}, \citet{kilian2009}, whose decomposition of oil price movements into supply, aggregate-demand and precautionary-demand components remains the organising framework, and \citet{baumeister2019}. On the current episode, \citet{dallasfed2026} map a range of Hormuz-closure scenarios into US headline and core PCE inflation, while \citet{dallasfed2026b} put the 2026 shortfall in global oil supply at more than twice the peak disruption of the 1973 crisis---the largest on record before 2026---and estimate that the US real GDP growth response today is only around one-sixth of the decline in the rest of the world---which implies a substantially larger effect on an economy like the UK's, and for which \emph{the UK counterpart is unwritten}. \citet{imf2026ports} add a transmission channel specific to chokepoint closure, showing with vessel-level AIS data that a 100-hour increase in port-to-port shipping time raises inflation by around 0.5 percentage points at a five-month peak. The aggregate scenario backdrop is \citet{weo2026} and \citet{worldbank2026}.

\subsection{Positioning}

Three gaps follow. First, the international consensus against price subsidies has never been tested against the peer-reviewed result that the optimum contains one, in the country where that result was estimated and on the policy choice now being made. Second, the pounds-versus-per-cent tension has never been displayed simultaneously on the same British households, and the within-decile uncompensated-loser count that \citet{sallee2019}, \citet{cronin2019} and \citet{douenne2020} all identify as decisive has never been reported for the UK proposals. Third, and most concretely, the Seventh Carbon Budget established ``macroeconomic model plus distributional annex'' as the expected format for UK official analysis---\citet{youngman2026} cite PolicyEngine in exactly that spirit, as a methodological precedent for wealth imputation rather than as an integrated component---and yet no one has built the macro-to-PolicyEngine link that the format presupposes. That link is what this paper constructs.
